\documentclass[11pt]{article}
\usepackage[localtheorems]{raeez-math-template}

\newtheorem{theorem}{Theorem}
\newtheorem{proposition}[theorem]{Proposition}
\newtheorem{conjecture}[theorem]{Conjecture}
\newtheorem{warning}[theorem]{Warning}

\providecommand{\kch}{\kappa_{\mathrm{ch}}}
\providecommand{\kcat}{\kappa_{\mathrm{cat}}}
\providecommand{\kBKM}{\kappa_{\mathrm{BKM}}}
\providecommand{\kfib}{\kappa_{\mathrm{fiber}}}
\providecommand{\fg}{\mathfrak{g}}
\providecommand{\ZZ}{\mathbb{Z}}
\providecommand{\CC}{\mathbb{C}}
\providecommand{\Sp}{\mathrm{Sp}}
\providecommand{\Aut}{\mathrm{Aut}}
\providecommand{\wt}{\operatorname{wt}}
\providecommand{\cO}{\mathcal{O}}
\providecommand{\Delf}{\Delta_5}
\providecommand{\PhiTen}{\Phi_{10}}
\providecommand{\Mtf}{M_{24}}
\providecommand{\phiK}{\phi^{K3}_{0,1}}

\title{Moonshine Atlases, K3 Twining, and Borcherds Constants}
\author{Raeez Lorgat}
\date{2026-04-22}

\begin{document}
\maketitle

\begin{abstract}
The K3 moonshine, umbral, Enriques, and Borcherds data use different
indexing sets.  The K3 elliptic genus twining system is indexed by the
26 conjugacy classes of \(M_{24}\).  Umbral moonshine is indexed by the
23 Niemeier root systems, with the Leech case giving Conway moonshine.
The Enriques denominator uses the Enriques lattice and its own
multiplier data.  Borcherds denominators require a Jacobi input,
multiplier, lattice, Weyl chamber, and root datum.  The automatic
numerical rule in the denominator lane is
\[
  \kBKM(\Phi)=\wt(\Phi)=\frac{c(0)}2
\]
for the chosen Borcherds input.
\end{abstract}

\section{The Four Atlases}

\begin{enumerate}[label=(\Alph*)]
\item \textbf{K3 elliptic-genus twining.}
For each conjugacy class \([g]\subset \Mtf\), Gaberdiel-Hohenegger-Volpato
attach a weak Jacobi form
\[
  \phi_g(\tau,z)
  =\operatorname{Tr}_{\mathcal H_{RR}}
  \left(g(-1)^{F+\widetilde F}q^{L_0-c/24}
  \bar q^{\widetilde L_0-\widetilde c/24}y^{J_0}\right)
  \in J^{w,\nu_g}_{0,1}(\Gamma_0(|g|)).
\]
The specialisation \(\phi_g(\tau,0)\) is the \(24\)-dimensional
Mathieu character \(\chi_{24}(g)\).  This datum is a one-variable
twining genus; denominator and twisted-twined interpretations require
additional structures.

\item \textbf{Umbral moonshine.}
Cheng-Duncan-Harvey attach mock modular forms to the 23 Niemeier
root systems through the quotient
\(\Aut(\Lambda)/W(\Lambda)\).  The lattice \(\Lambda=24A_1\) gives the
Mathieu case.  The remaining 22 root systems are separate umbral
systems, and the rootless Leech lattice belongs to the Conway branch.
The words ``Mathieu'' and ``umbral'' name different indexing rules:
conjugacy classes of \(M_{24}\) in the first case and Niemeier root
systems in the second.

\item \textbf{Enriques denominator data.}
An Enriques surface has \(K_S\) nontrivial
2-torsion and \(\chi(\cO_S)=1\), so it lies outside the Calabi-Yau
surface class.  Its K3 cover is governed by a
fixed-point-free anti-symplectic involution.  The symplectic Mathieu
class \(2A\) has Frame shape \(1^8 2^8\) and K3 fixed-locus Euler
number \(8\); it is a different row.  The Enriques automorphic product
uses the Enriques lattice and its multiplier.

\item \textbf{Borcherds denominators.}
A Borcherds denominator requires a weak Jacobi input, a group or cover,
a character or multiplier, a lattice of signature \((2,n)\), a Weyl
chamber, and a root-multiplicity convention.  The product weight is
\(c(0)/2\) in that convention.  A Mathieu character-table row supplies
the conjugacy-class datum only.
\end{enumerate}

\section{Numerical Normalisations}

\subsection{Ordinary \(M_{24}\) Frame Shapes}

\begin{proposition}[The 26 ordinary Mathieu rows]
\label{prop:m24-ordinary-26-rows}
In the 24-letter permutation normalisation, the ordinary Mathieu
twining table has the following 26 rows:
\begin{center}
\small
\begin{tabular}{c|c|c|c}
class & $|g|$ & Frame shape & $\chi_{24}(g)$\\
\hline
1A & 1 & $1^{24}$ & 24\\
2A & 2 & $1^8 2^8$ & 8\\
2B & 2 & $2^{12}$ & 0\\
3A & 3 & $1^6 3^6$ & 6\\
3B & 3 & $3^8$ & 0\\
4A & 4 & $2^4 4^4$ & 0\\
4B & 4 & $1^4 2^2 4^4$ & 4\\
4C & 4 & $4^6$ & 0\\
5A & 5 & $1^4 5^4$ & 4\\
6A & 6 & $1^2 2^2 3^2 6^2$ & 2\\
6B & 6 & $6^4$ & 0\\
7A & 7 & $1^3 7^3$ & 3\\
7B & 7 & $1^3 7^3$ & 3\\
8A & 8 & $1^2 2^1 4^1 8^2$ & 2\\
10A & 10 & $2^2 10^2$ & 0\\
11A & 11 & $1^2 11^2$ & 2\\
12A & 12 & $2^1 4^1 6^1 12^1$ & 0\\
12B & 12 & $12^2$ & 0\\
14A & 14 & $1^1 2^1 7^1 14^1$ & 1\\
14B & 14 & $1^1 2^1 7^1 14^1$ & 1\\
15A & 15 & $1^1 3^1 5^1 15^1$ & 1\\
15B & 15 & $1^1 3^1 5^1 15^1$ & 1\\
21A & 21 & $3^1 21^1$ & 0\\
21B & 21 & $3^1 21^1$ & 0\\
23A & 23 & $1^1 23^1$ & 1\\
23B & 23 & $1^1 23^1$ & 1
\end{tabular}
\end{center}
The row \(3B\) is part of the ordinary table.  The corrected rows
\[
4A=2^4 4^4,\quad
4B=1^4 2^2 4^4,\quad
6B=6^4,\quad
10A=2^2 10^2,\quad
12A=2^1 4^1 6^1 12^1,\quad
12B=12^2
\]
are the ones used by the local Mathieu compute oracle.
\end{proposition}

\begin{proof}
Each Frame shape \(\prod_i i^{a_i}\) is the cycle type of \(g\) in the
natural 24-letter permutation representation, hence
\(\sum_i i a_i=24\), and the displayed character value is \(a_1\).
The local source
\(\texttt{compute/lib/mathieu\_moonshine\_yangian.py}\) stores exactly
these 26 rows and checks both identities by
\(\texttt{verify\_frame\_shapes()}\).
\end{proof}

\subsection{The Primitive Level-One Denominator}

The index-one generator is normalised by
\[
  \phi_{0,1}(\tau,z)=y^{-1}+10+y+O(q).
\]
Hence the discriminant-zero coefficient is \(c_1(0)=10\).  Borcherds'
weight theorem gives
\[
  \kBKM(\Delf)=\wt(\Delf)=10/2=5.
\]
The Igusa form \(\PhiTen=\Delf^2\) has weight \(10\).  The square is the
scalar genus-two partition function normalisation; the primitive BKM
denominator is \(\Delf\) of weight \(5\).

\subsection{CHL-Averaged Denominators}

For the cyclic CHL denominator inputs used in the \(K3\times E\)
Borcherds lane, the local compute witnesses
\texttt{compute/lib/diagonal\_siegel\_cy\_orbifolds.py} and
the universal BKM-weight witness use the row
\[
\begin{array}{c|c|c|c}
N & \text{frame shape} & c_N^{\mathrm{CHL}}(0) &
\kBKM(\Phi_N^{\mathrm{CHL}})\\
\hline
1 & 1^{24} & 10 & 5\\
2 & 1^{8}2^{8} & 8 & 4\\
3 & 1^{6}3^{6} & 6 & 3\\
4 & 1^{4}2^{2}4^{4} & 4 & 2\\
6 & 1^{2}2^{2}3^{2}6^{2} & 2 & 1
\end{array}
\]
This is an averaged orbifold denominator row.  Ordinary \(M_{24}\)
twining coefficients live in Proposition~\ref{prop:m24-ordinary-26-rows}
and in the Jacobi forms \(\phi_g\).

\subsection{Singly Twined Mathieu Inputs}

For the ten low-order \(M_{24}\) classes
\[
1A,\ 2A,\ 2B,\ 3A,\ 3B,\ 4A,\ 4B,\ 4C,\ 6A,\ 6B
\]
the \(q^0y^0\)-coefficients of the ordinary twined K3 elliptic genera
are
\[
  c^{(g)}(0,0)=(20,4,-4,2,-4,0,4,-4,0,0).
\]
The half-vector
\[
  c^{(g)}(0,0)/2=(10,2,-2,1,-2,0,2,-2,0,0)
\]
is a singly twined scalar additive-lift vector.  The CHL denominator
row is \((5,4,3,2,1)\).  Rows with zero or negative half-coefficient
are ordinary twining data; primitive positive-weight BKM denominators
use Borcherds input constants.

\subsection{Cl\'ery-Gritsenko Diagonal-Divisor Forms}

The Cl\'ery-Gritsenko catalogue is indexed by triples \((t,N;k)\).  The
CHL order alone is insufficient:
\[
\begin{array}{c|c|c|c}
\text{form} & (t,N) & \text{group} & \wt\\
\hline
\Delf & (1,1) & \Sp_4(\ZZ) & 5\\
\Delta_2 & (2,1) & \Gamma_2 & 2\\
\nabla_3 & (1,2) & \Gamma_0^{(2)}(2) & 3\\
\Delta_1 & (3,1) & \Gamma_3 & 1\\
\nabla_2 & (1,3) & \Gamma_0^{(2)}(3) & 2\\
\Delta_{1/2} & (4,1) & \Gamma_4 & 1/2\\
\nabla_{3/2} & (1,4) & \Gamma_0^{(2)}(4) & 3/2\\
Q_1 & (2,2) & \Gamma_2(2) & 1
\end{array}
\]
The \(\Delf\) row is the common intersection of the CHL ladder and the
Cl\'ery-Gritsenko catalogue.  Ordinary Mathieu twining supplies the
identity Jacobi input; the paramodular group, product, and BKM root data
are additional structure.  The half-integral rows require the stated
cover and multiplier.  A transfer to a CHL or Mathieu table requires an
explicit row map.

\section{K3, Enriques, and Persson-Volpato}

The ordinary twined genus \(\phi_g\) exists for every \(M_{24}\) class.
The twisted-twined genus \(\phi_{g,h}\) requires an actual commuting
pair \(gh=hg\), a projective action of \(C_{M_{24}}(g)\) on the
\(g\)-twisted sector, and the Persson-Volpato anomaly multiplier.  The
second-quantised product of the twisted-twined genus is a Siegel
modular object after the anomaly correction.  It is a BKM denominator
only after the lattice, chamber, Weyl vector, cover, and product
exponents have been supplied.

\begin{warning}[Class-row separation]
Canonical maps from the 26 \(M_{24}\) conjugacy classes to CHL or
Cl\'ery-Gritsenko denominator rows require extra data.  A valid row
assertion names:
the conjugacy class or commuting pair, the multiplier, the Jacobi input,
the paramodular group or orthogonal group, the product weight, and the
root datum.  With only the class row, the mathematical object is a K3
twining genus or a twisted-twined genus.
\end{warning}

\section{Consequences for Vol III}

\[
\begin{array}{c|c|c}
\text{object} & \text{indexing set} & \text{constant}\\
\hline
\Delf & \phi_{0,1} & \kBKM=5\\
\PhiTen=\Delf^2 & \text{Igusa square} & \wt=10\\
\Phi_N^{\mathrm{CHL}} & N=1,2,3,4,6 & \kBKM=(5,4,3,2,1)\\
\phi_g & [g]\subset M_{24} & \chi_{24}(g),\ c^{(g)}(0,0)\\
\text{umbral }H^{(\Lambda)} & 23\text{ Niemeier root systems} &
\text{mock modular shadow}\\
\text{Enriques product} & \text{Enriques lattice} &
\text{lattice and multiplier dependent}
\end{array}
\]

The four Vol III invariants remain separated:
\[
  \kcat(K3\times E)=0,\qquad
  \kch^{\mathrm{Heis}}(K3\times E)=3,\qquad
  \kBKM(\Delf)=5,\qquad
  \kfib(K3)=24.
\]
The equality \(3+2=5\) at the K3 fibre is a comparison of different
constructions at \(N=1\).  The denominator formula is
\(\kBKM(\Phi)=c(0)/2\) for the chosen Borcherds input.

\section{Source Anchors}

\begin{itemize}
\item Borcherds 1998, \emph{Automorphic forms with singularities on
Grassmannians}, Theorem 13.3: product weight \(c(0)/2\).
\item Gritsenko-Nikulin 1998 and Gritsenko 1999: \(\Delf\) and the
rank-three denominator construction.
\item Eguchi-Ooguri-Tachikawa 2011; Gaberdiel-Hohenegger-Volpato
2012: \(M_{24}\)-twined K3 elliptic genera.
\item Cheng-Duncan-Harvey 2014: umbral moonshine and the Niemeier
root-system indexing.
\item Persson-Volpato 2012 and Gaberdiel-Persson-Volpato 2013:
twisted-twined genera and anomaly phases.
\item Allcock 2000 and the Enriques-lattice Borcherds literature:
the Enriques denominator branch.
\end{itemize}

\end{document}
